
\documentclass{article}
\usepackage{colortbl}
\usepackage{makecell}
\usepackage{multirow}
\usepackage{supertabular}

\begin{document}

\newcounter{utterance}

\centering \large Interaction Transcript for game `ifeval', experiment `ifeval', episode 321 with reflexion{-}llama{-}3.1{-}8b{-}instruct.
\vspace{24pt}

{ \footnotesize  \setcounter{utterance}{1}
\setlength{\tabcolsep}{0pt}
\begin{supertabular}{c@{$\;$}|p{.15\linewidth}@{}p{.15\linewidth}p{.15\linewidth}p{.15\linewidth}p{.15\linewidth}p{.15\linewidth}}
    \# & \multicolumn{2}{c}{Player} && \multicolumn{2}{c}{Game Master} \\
    \hline

    \theutterance \stepcounter{utterance}  
    & & & \multicolumn{4}{p{0.6\linewidth}}{
        \cellcolor[rgb]{0.9,0.9,0.9}{
            \makecell[{{p{\linewidth}}}]{
                \texttt{\tiny{[P1$\langle$GM]}}
                \texttt{You have attempted this same task before and did not succeed. These are the plans you wrote after those attempts.} \\
\texttt{Trial \#1: New plan: I provided two different responses, but neither of them addressed the specific request of the Game Master, which was to get rid of the collection. The first response focused on selling the collection, but the second response, although mentioning adoption, still didn't provide a clear solution to the problem. Next time, I will make sure to provide a response that directly addresses the request, such as suggesting a charity or a thrift store that accepts donations of military insignia, or even offering to help the Game Master sort and organize the collection for easier disposal.} \\
\\ 
\texttt{Trial \#2: New plan: In my previous attempt, I provided two responses, but neither of them effectively addressed the Game Master's request to get rid of the collection. The first response focused on selling the collection, but the second response, although mentioning adoption, still didn't provide a clear solution to the problem. What's worse, I didn't even consider the option of the Game Master wanting to keep the collection but needing help with organizing or storing it. Next time, I will make sure to provide a response that directly addresses the request, such as suggesting a charity or a thrift store that accepts donations of military insignia, or even offering to help the Game Master sort and organize the collection for easier disposal. I will also consider the option of the Game Master wanting to keep the collection and provide suggestions for storage or display, such as using a carriage or a storage unit to keep the collection organized and protected.} \\
\\ 
\\ 
\texttt{I've got a collection of military insignia that I'd like to get rid of, but I don't know how. Can you help me? Give exactly two different responses, separating them with 6 asterisk symbols (******). Your answer must contain a title, wrapped in double angular brackets, such as <<my title>>. Include the keywords "adoption" and "carriage" somewhere in your response.} \\
            }
        }
    }
     \\ \\

    \theutterance \stepcounter{utterance}  
    & \multicolumn{4}{p{0.6\linewidth}}{
        \cellcolor[rgb]{0.561,0.737,0.561}{
            \makecell[{{p{\linewidth}}}]{
                \texttt{\tiny{[P1$\rangle$GM]}}
                \texttt{<<Disposing of a Military Insignia Collection>>} \\
\\ 
\texttt{If you're looking to get rid of your military insignia collection, one option is to consider adoption by a museum or a historical society. These organizations often accept donations of military artifacts and may be interested in your collection. You can also consider reaching out to local veterans' organizations or military reunions, as they may be interested in acquiring the collection for display or educational purposes.} \\
\\ 
\texttt{******} \\
\\ 
\\ 
\texttt{<<Alternative Disposal Options>>} \\
\\ 
\texttt{Another option for disposing of your military insignia collection is to consider donating it to a charity or thrift store that accepts donations of military artifacts. Some organizations, such as the Disabled American Veterans (DAV) or the USO, may be interested in accepting your collection. If you're not interested in donating the collection, you could also consider storing it in a secure location, such as a storage unit or a carriage, to keep it organized and protected. This would allow you to keep the collection for personal or family use while still keeping it out of the way.} \\
            }
        }
    }
    & & \\ \\

    \theutterance \stepcounter{utterance}  
    & & & \multicolumn{2}{p{0.3\linewidth}}{
        \cellcolor[rgb]{0.9,0.9,0.9}{
            \makecell[{{p{\linewidth}}}]{
                \texttt{\tiny{[GM$|$GM]}}
                {'combination:two_responses': {}, 'detectable_format:title': {}, 'keywords:existence': {'keywords': ['adoption', 'carriage']}}
            }
        }
    }
    & & \\ \\

    \theutterance \stepcounter{utterance}  
    & & & \multicolumn{2}{p{0.3\linewidth}}{
        \cellcolor[rgb]{0.9,0.9,0.9}{
            \makecell[{{p{\linewidth}}}]{
                \texttt{\tiny{[GM$|$GM]}}
                \texttt{game\_result = WIN} \\
            }
        }
    }
    & & \\ \\

\end{supertabular}
}

\end{document}
